% ═══════════════════════════════════════════════════════════════════════════════
% CASO 3: GUARDERÍA ABC - INSTITUTO ATENEO DE COLIMA
% Estudio Documental: Corrupción en México (2006-2012)
% ═══════════════════════════════════════════════════════════════════════════════

\nota{
  Este capítulo analiza la desregulación y conflicto de interés en el caso 
  Guardería ABC mediante evidencia de CNDH (2009), SCJN (2010, 2025) y 
  doctrina de Serna de la Garza (2010) y Osorno (2010).
}

\section{Retiro estatal, expansión de la subrogación y conflicto de interés (1995--2009)}

El incendio de la Guardería ABC en Hermosillo, Sonora, ocurrido el 5 de junio de 2009 ---en el que fallecieron 49 niñas y niños y más de 100 menores sufrieron lesiones graves---, evidenció los riesgos extremos advertidos por Rose-Ackerman (1999), Rezzoagli et al. (2021), Sandoval Ballesteros (2009) y Serna de la Garza (2010) en la interfaz entre el Estado regulador y los operadores privados cuando confluyen la desregulación, el conflicto de interés y la ausencia de vigilancia independiente.

El seguro de guarderías del Instituto Mexicano del Seguro Social (IMSS) surgió en 1973 como un servicio público directo financiado con cuotas obrero-patronales. A partir de 1995 se impulsó el esquema «vecinal-comunitario» de subrogación a particulares: entre 2000 y 2009, las estancias subrogadas aumentaron de 757 a 1,374 ---un crecimiento de 81.5\,\%---, mientras que las operadas directamente por el IMSS permanecieron congeladas en 142 \citep{leal2010}.

% FIGURA PENDIENTE
\begin{figure}[H]
\centering
\includegraphics[width=0.9\textwidth]{img/image2.png} % Figura 2: Expansión guarderías IMSS 2000-2009
\caption{Expansión de guarderías subrogadas frente a guarderías de operación directa del IMSS (2000--2009). Fuente: Elaboración propia con datos de Leal (2010).}
\label{fig:guarderias-imss}
\end{figure}

La concesión de la Guardería ABC fue otorgada a una sociedad civil integrada por particulares con vínculos familiares y políticos con altos funcionarios estatales y federales ---entre quienes figuraban Marcia Matilde Altagracia Gómez del Campo Tonella (prima de la entonces primera dama Margarita Zavala) y Sandra Lucía Téllez Nieves, además de socios ligados al gobierno de Sonora--- \citep{scjn2009,osorno2010}. Como documentaron la investigación periodística de Osorno (2010) y los expedientes de derechos humanos, la estancia operaba en una nave industrial adaptada con techumbre inflamable, salidas de emergencia deficientes y únicamente 14 de los 40 trabajadores reglamentarios el día del siniestro, colindando pared con pared con una bodega de archivo y vehículos de la Secretaría de Hacienda estatal \citep{leal2010,nunez2026}.

La Recomendación 49/2009 de la CNDH acreditó que el siniestro era previsible y evitable: existía un dictamen arquitectónico fechado en julio de 2005 que ya había señalado riesgos graves en el inmueble sin que el IMSS ni Protección Civil ordenaran correcciones o la suspensión del servicio \citep{cndh2009}. En concordancia con Serna de la Garza (2010) y Arellano Gault (2012), la transferencia desregulada del servicio y las redes informales de cercanía política diluyeron la responsabilidad pública y neutralizaron los controles preventivos.

\section{Resoluciones constitucionales e internacionales (2010--2026)}

La trayectoria jurídica del expediente consolidó precedentes relevantes en el sistema constitucional mexicano y en el ámbito interamericano:

\begin{itemize}
  \item \textbf{Suprema Corte de Justicia de la Nación (SCJN):} En el dictamen final de la Facultad de Investigación 1/2009 (ponente ministro Arturo Zaldívar) y posteriormente en la sentencia del Amparo Directo 14/2019, el alto tribunal acreditó la existencia de un «desorden generalizado» en el sistema nacional de guarderías y una subrogación efectuada al margen del artículo 213 de la Ley del Seguro Social, aunque en 2010 no se alcanzó mayoría para fincar responsabilidad directa a los titulares nacionales del IMSS \citep{scjn2010,scjn2025,zaldivar2010}.
  
  \item \textbf{Comisión Interamericana de Derechos Humanos (CIDH):} En julio de 2020, mediante el Informe No. 185/20, admitió la petición internacional al constatar la omisión de fiscalizaciones de oficio por parte del IMSS y la ampliación contractual pese a los peligros evidentes de la bodega contigua \citep{cidh2020}.
  
  \item \textbf{Declaratoria de imprescriptibilidad penal (2026):} El 19 de junio de 2026, el Pleno de la SCJN declaró imprescriptible la acción penal en el caso Guardería ABC, determinando que el paso del tiempo no extingue el derecho a la justicia ni a la reparación integral ante violaciones graves a los derechos de la niñez \citep{nunez2026}.
\end{itemize}

\section{Casos de contraste contemporáneo: daño real en derechos humanos frente a encuadre mediático (framing)}

Con el fin de responder la cuarta pregunta de investigación sin alterar la delimitación nuclear 2006--2012, se examinaron dos casos de contraste contemporáneo que permitieron contrastar el criterio de afectación verificable a derechos humanos \citep{cidh2019,acnudh} frente a las distorsiones del encuadre mediático \citep{solamorales2019}:

\subsection{Contraste 1 --- Sector salud en Veracruz}

En enero de 2017 se viralizó la acusación de que durante la gubernatura de Javier Duarte se había aplicado agua destilada a niños con cáncer. Tras veinte visitas de inspección sanitaria, la Secretaría de Salud federal y la COFEPRIS informaron oficialmente que no se encontró evidencia de medicamentos falsos ni de agua destilada suministrada a pacientes oncológicos \citep{cofepris2017}. No obstante, la narrativa viral eclipsó parcialmente un saqueo masivo plenamente comprobado: la propia COFEPRIS aseguró 46,984 pruebas rápidas de VIH caducas y sin registro y 16.8 toneladas de medicamentos caducos, mientras auditorías e investigaciones documentaron 2,933 millones de pesos desviados en obra hospitalaria (183 obras inconclusas y 12 hospitales abandonados), más de 3,617 millones de pesos transferidos a empresas fantasma y 276.6 millones adicionales en factureras ---hechos que derivaron en nuevas imputaciones de la FGR en diciembre de 2025--- \citep{connectas,animalpolitico,santiago2025,angel2025}.

\subsection{Contraste 2 --- Megainfraestructura energética (Dos Bocas)}

De modo comparable al símbolo reduccionista de «la barda» en la Refinería Bicentenario de Tula, el debate público sobre la Refinería Olmeca (Dos Bocas) tendió a polarizarse en consignas sobre si la planta «no refina un solo litro». Los reportes técnicos oficiales mostraron una realidad más matizada pero financieramente crítica: el costo se elevó de 8,000 millones de dólares proyectados a más de 21,000 millones al cierre de 2024 (\textasciitilde2.6 veces el presupuesto original) y, aunque la producción alcanzó 212,000--263,000 barriles diarios hacia finales de 2025 (frente a su capacidad nominal de 340,000), descendió a 141,000 barriles diarios en junio de 2026 tras un incendio en la planta de coque y una falla eléctrica \citep{alegria2026,martinezescamilla2026,celaya2026}. Reducir estos proyectos a eslóganes partidistas desplazó la fiscalización rigurosa sobre adjudicaciones, sobrecostos y costo de oportunidad fiscal.

\section{Ejes de corrupción en México jerarquizados por alcance poblacional}

Al jerarquizar las modalidades de corrupción documentadas en México conforme al criterio objetivo de número de afectados y vulneración de derechos fundamentales ---al margen de afinidades partidistas---, emergieron cinco ejes de impacto estructural:

\begin{table}[H]
\centering
\caption{Ejes de corrupción en México jerarquizados por alcance poblacional y derechos vulnerados}
\label{tab:ejes_corrupcion}
\small
\begin{tabular}{@{}p{3.5cm}p{4cm}p{5.5cm}@{}}
\toprule
\textbf{Eje de corrupción} & \textbf{Casos / evidencia emblemática} & \textbf{Magnitud verificable del daño y población afectada} \\
\midrule
Captura de seguridad pública y macrocriminalidad & Genaro García Luna (SSP, 2006--2012) y colusión local con crimen organizado. & Afectación sobre millones de habitantes por violencia, narcotráfico y extorsión. Condena penal de >38 años (2024) y civiles por \textasciitilde2,500 MDD (2025) y >578 MDD (2026). \\
\addlinespace
Desfalco a programas sociales y abasto alimentario & Segalmex (2019--2024) y La Estafa Maestra (2012--2018). & Irregularidades en Segalmex por >15,000--15,300 mdp destinados a pequeños productores y abasto rural; desvíos en La Estafa Maestra por >7,000 mdp en programas sociales \citep{perez2024,mccitablero}. \\
\addlinespace
Desvíos y deterioro del sector salud & Saqueo hospitalario estatal (Veracruz) e irregularidades en transferencias federales. & Impacto directo en millones de usuarios sin seguro privado; en noviembre de 2026 la ASF aún señaló 285.8 mdp por aclarar y 553.3 mdp reintegrados en salud estatal \citep{ramos2026}; IMSS con 43.3\,\% de satisfacción \citep{inegi2026}. \\
\addlinespace
Desregulación con conflicto de interés letal & Subrogación irregular en la Guardería ABC (2009). & 49 niñas y niños fallecidos y más de 100 lesionados graves; «desorden generalizado» acreditado por la SCJN e imprescriptibilidad penal decretada en junio de 2026. \\
\addlinespace
Corrupción cotidiana («mordida» en trámites) & Medición nacional ENCIG 2025 (INEGI, publ. mayo 2026). & Impuesto regresivo directo: afectó a 15.6\,\% de quienes gestionaron trámites y extrajo \$17,707.4 millones de pesos anuales (\$3,865 pesos promedio por víctima). \\
\bottomrule
\end{tabular}
\end{table}